\definicion{Validación W3C}{Proceso de comprobación del código (HTML, CSS) contra los estándares definidos por el World Wide Web Consortium, asegurando la validez y la consistencia entre navegadores.}
\subtituloA{Pasos recomendados para un portal web}

\lista{
  \textbf{Paso 1:} Codificar los códigos HTML, JavaScript y CSS en el documento. |
  \textbf{Paso 2:} Guardar con extensiones correctas y ejecutar en navegadores; las modificaciones se ven recargando. |
  \textbf{Paso 3:} Validar el código utilizando el validador oficial de W3C (http://validator.w3.org/).
}
\nota{Para el examen}{
La estandarización W3C es clave para garantizar la portabilidad y el mejor resultado entre diferentes navegadores, no solo para la funcionalidad local.
}
